\documentclass[aspectratio=169]{beamer}
\usetheme{Madrid}
\usecolortheme{whale}
\usepackage[utf8]{inputenc}
\usepackage[spanish]{babel}
\usepackage{amsmath,amssymb}
\usepackage{graphicx}
\usepackage{booktabs}
\usepackage{listings}
\usepackage{xcolor}

\definecolor{codegreen}{rgb}{0,0.6,0}
\definecolor{codegray}{rgb}{0.5,0.5,0.5}
\definecolor{codepurple}{rgb}{0.58,0,0.82}
\definecolor{backcolour}{rgb}{0.95,0.95,0.92}

\lstdefinestyle{pythonstyle}{
    backgroundcolor=\color{backcolour},
    commentstyle=\color{codegreen},
    keywordstyle=\color{blue},
    numberstyle=\tiny\color{codegray},
    stringstyle=\color{codepurple},
    basicstyle=\ttfamily\footnotesize,
    breakatwhitespace=false,
    breaklines=true,
    captionpos=b,
    keepspaces=true,
    showspaces=false,
    showstringspaces=false,
    showtabs=false,
    tabsize=2
}

\lstset{style=pythonstyle}

\title{Selección de Variables y Regresión Polinomial}
\subtitle{Elegir el mejor modelo}
\author{Métodos de Análisis de Datos 1}
\institute{Universidad Nacional del Sur}
\date{Unidad 3 - Clase 18}

\begin{document}

\frame{\titlepage}

\begin{frame}{Contenidos}
\tableofcontents
\end{frame}

\section{Selección de Variables}

\begin{frame}{¿Por qué seleccionar variables?}
Con muchas variables candidatas, ¿cuáles incluir?

\vspace{0.3cm}
\textbf{Tensión:}
\begin{itemize}
    \item Más variables $\Rightarrow$ mejor ajuste ($R^2$ más alto)
    \item Pero: sobreajuste, interpretación difícil, multicolinealidad
\end{itemize}

\vspace{0.5cm}
\textbf{Ventajas de modelos simples:}
\begin{enumerate}
    \item Más interpretables
    \item Mejor generalización (menor sobreajuste)
    \item Menor costo computacional
    \item Predicciones más estables
\end{enumerate}

\vspace{0.3cm}
\textbf{Objetivo:} Encontrar el modelo más simple que explique bien los datos
\end{frame}

\begin{frame}{Criterios de selección}
\textbf{AIC (Criterio de Información de Akaike):}
\[
\text{AIC} = n \ln\left(\frac{\text{SCR}}{n}\right) + 2(p+1)
\]

\textbf{BIC (Criterio Bayesiano):}
\[
\text{BIC} = n \ln\left(\frac{\text{SCR}}{n}\right) + \ln(n)(p+1)
\]

\vspace{0.5cm}
\textbf{Interpretación:} Menor es mejor

\vspace{0.3cm}
\textbf{Diferencia:} BIC penaliza más la complejidad que AIC
\begin{itemize}
    \item AIC $\Rightarrow$ modelos más complejos
    \item BIC $\Rightarrow$ modelos más simples (preferido para $n$ grande)
\end{itemize}
\end{frame}

\begin{frame}{Métodos de búsqueda}
\textbf{1. Best Subset (mejor subconjunto):}
\begin{itemize}
    \item Probar todos los modelos posibles: $2^p$
    \item Elegir el de menor AIC/BIC
    \item Computacionalmente costoso para $p > 15$
\end{itemize}

\vspace{0.3cm}
\textbf{2. Forward (hacia adelante):}
\begin{itemize}
    \item Empezar con modelo vacío
    \item Agregar la variable que más mejora el criterio
    \item Repetir hasta que ninguna mejore
\end{itemize}

\vspace{0.3cm}
\textbf{3. Backward (hacia atrás):}
\begin{itemize}
    \item Empezar con modelo completo
    \item Eliminar la variable menos significativa
    \item Repetir hasta que todas sean significativas
\end{itemize}
\end{frame}

\begin{frame}{Método Stepwise}
\textbf{Stepwise (por pasos):} Combina forward y backward

\vspace{0.3cm}
\textbf{Algoritmo:}
\begin{enumerate}
    \item Empezar con modelo vacío (o completo)
    \item En cada paso:
    \begin{itemize}
        \item Considerar agregar una variable (forward)
        \item Considerar eliminar una variable (backward)
    \end{itemize}
    \item Elegir la acción que más mejora el criterio
    \item Repetir hasta convergencia
\end{enumerate}

\vspace{0.5cm}
\begin{alertblock}{Advertencia}
Los métodos automáticos son heurísticos. No garantizan encontrar el mejor modelo. Siempre combinar con conocimiento del dominio.
\end{alertblock}
\end{frame}

\begin{frame}{Validación cruzada}
\textbf{Problema:} AIC/BIC se calculan en los mismos datos usados para ajustar

\vspace{0.3cm}
\textbf{Solución:} Validación cruzada

\vspace{0.3cm}
\textbf{K-fold CV:}
\begin{enumerate}
    \item Dividir datos en $K$ partes (folds)
    \item Para cada fold:
    \begin{itemize}
        \item Ajustar modelo en $K-1$ folds
        \item Predecir en el fold restante
        \item Calcular error
    \end{itemize}
    \item Error CV = promedio de los $K$ errores
\end{enumerate}

\vspace{0.3cm}
\textbf{Común:} $K = 5$ o $K = 10$

\vspace{0.3cm}
\textbf{Ventaja:} Estima error de predicción en datos nuevos
\end{frame}

\section{Regresión Polinomial}

\begin{frame}{¿Cuándo usar polinomios?}
Cuando la relación entre $X$ e $Y$ no es lineal.

\vspace{0.3cm}
\textbf{Modelo polinomial de grado $d$:}
\[
Y = \beta_0 + \beta_1 X + \beta_2 X^2 + \cdots + \beta_d X^d + \epsilon
\]

\vspace{0.5cm}
\textbf{Nota importante:} Aunque la relación con $X$ no es lineal, el modelo es \textbf{lineal en los parámetros} $\beta_j$. Por lo tanto:
\begin{itemize}
    \item Se puede usar MCO
    \item Todas las técnicas de regresión múltiple aplican
\end{itemize}

\vspace{0.3cm}
\textbf{Implementación:} Crear variables $X_2 = X^2$, $X_3 = X^3$, etc., y ajustar regresión múltiple estándar.
\end{frame}

\begin{frame}{Grados comunes}
\textbf{Grado 2 (cuadrática):}
\[
Y = \beta_0 + \beta_1 X + \beta_2 X^2 + \epsilon
\]
\begin{itemize}
    \item Forma de U o U invertida
    \item Un punto de inflexión
\end{itemize}

\vspace{0.3cm}
\textbf{Grado 3 (cúbica):}
\[
Y = \beta_0 + \beta_1 X + \beta_2 X^2 + \beta_3 X^3 + \epsilon
\]
\begin{itemize}
    \item Forma de S
    \item Dos puntos de inflexión
\end{itemize}

\vspace{0.5cm}
\begin{alertblock}{Cuidado}
Grados muy altos ($d > 4$) tienden a sobreajustar
\end{alertblock}
\end{frame}

\begin{frame}{Selección del grado}
¿Qué grado $d$ elegir?

\vspace{0.3cm}
\textbf{Métodos:}
\begin{enumerate}
    \item \textbf{Visualización:} Graficar ajustes para distintos $d$
    \item \textbf{AIC/BIC:} Elegir $d$ con menor AIC/BIC
    \item \textbf{Validación cruzada:} Elegir $d$ con menor error CV
    \item \textbf{Test F jerárquico:} Comparar modelos anidados
\end{enumerate}

\vspace{0.5cm}
\textbf{Recomendación:} Empezar simple ($d=2$) y aumentar solo si es necesario
\end{frame}

\begin{frame}{Interpretación de polinomios}
\textbf{Problema:} Los coeficientes individuales son difíciles de interpretar

\vspace{0.3cm}
\textbf{Ejemplo:} $\hat{Y} = 10 + 3X - 0.5X^2$
\begin{itemize}
    \item $\beta_1 = 3$ no significa "por cada unidad de $X$, $Y$ aumenta 3"
    \item El efecto de $X$ depende del valor de $X$ (por el término $X^2$)
\end{itemize}

\vspace{0.5cm}
\textbf{Mejor enfoque:}
\begin{itemize}
    \item Graficar la curva ajustada
    \item Calcular efectos marginales para valores específicos de $X$
    \item Identificar máximos/mínimos
\end{itemize}
\end{frame}

\begin{frame}{Máximos y mínimos}
Para una cuadrática $Y = \beta_0 + \beta_1 X + \beta_2 X^2$:

\vspace{0.3cm}
\textbf{Punto crítico:}
\[
\frac{dY}{dX} = \beta_1 + 2\beta_2 X = 0 \quad \Rightarrow \quad X^* = -\frac{\beta_1}{2\beta_2}
\]

\vspace{0.3cm}
\begin{itemize}
    \item Si $\beta_2 < 0$ $\Rightarrow$ $X^*$ es un máximo
    \item Si $\beta_2 > 0$ $\Rightarrow$ $X^*$ es un mínimo
\end{itemize}

\vspace{0.5cm}
\textbf{Ejemplo:} $\hat{Y} = 10 + 3X - 0.5X^2$
\[
X^* = -\frac{3}{2(-0.5)} = 3
\]
El precio máximo se alcanza en $X = 3$
\end{frame}

\section{Implementación}

\begin{frame}[fragile]{Python: Selección stepwise}
\begin{lstlisting}[language=Python]
import statsmodels.api as sm

# No hay stepwise nativo, pero se puede hacer manual
# O usar libreria mlxtend
from mlxtend.feature_selection import SequentialFeatureSelector as SFS
from sklearn.linear_model import LinearRegression

# Forward selection
sfs = SFS(LinearRegression(), 
          k_features='best',
          forward=True, 
          scoring='neg_mean_squared_error',
          cv=5)

sfs.fit(X, y)
print(f"Variables seleccionadas: {sfs.k_feature_names_}")

# O calcular AIC/BIC manualmente
print(f"AIC: {modelo.aic:.2f}")
print(f"BIC: {modelo.bic:.2f}")
\end{lstlisting}
\end{frame}

\begin{frame}[fragile]{Python: Regresión polinomial}
\begin{lstlisting}[language=Python]
import numpy as np
from sklearn.preprocessing import PolynomialFeatures

# Crear features polinomiales
poly = PolynomialFeatures(degree=2, include_bias=False)
X_poly = poly.fit_transform(X)

# Ajustar modelo
modelo = LinearRegression().fit(X_poly, y)

# O con statsmodels y formula
from statsmodels.formula.api import ols
modelo = ols('Y ~ X + I(X**2)', data=data).fit()

# Para grado 3
modelo = ols('Y ~ X + I(X**2) + I(X**3)', data=data).fit()
\end{lstlisting}
\end{frame}

\begin{frame}[fragile]{R: Selección stepwise}
\begin{lstlisting}[language=R]
# Modelo completo
modelo_completo <- lm(Y ~ ., data = datos)

# Stepwise (ambas direcciones)
modelo_step <- step(modelo_completo, direction = "both")

# Forward
modelo_forward <- step(lm(Y ~ 1, data = datos),
                       scope = list(lower = ~1, 
                                    upper = modelo_completo),
                       direction = "forward")

# Backward
modelo_backward <- step(modelo_completo, 
                        direction = "backward")

# Ver criterios
AIC(modelo_step)
BIC(modelo_step)
\end{lstlisting}
\end{frame}

\begin{frame}[fragile]{R: Regresión polinomial}
\begin{lstlisting}[language=R]
# Grado 2
modelo_poly2 <- lm(Y ~ X + I(X^2), data = datos)

# Grado 3
modelo_poly3 <- lm(Y ~ X + I(X^2) + I(X^3), data = datos)

# O usar poly() que crea polinomios ortogonales
modelo_ortho <- lm(Y ~ poly(X, degree = 3), data = datos)

# Graficar
x_seq <- seq(min(datos$X), max(datos$X), length.out = 100)
y_pred <- predict(modelo_poly2, 
                  newdata = data.frame(X = x_seq))

plot(datos$X, datos$Y)
lines(x_seq, y_pred, col = "red", lwd = 2)
\end{lstlisting}
\end{frame}

\section{Resumen}

\begin{frame}{Resumen}
\textbf{Selección de variables:}
\begin{itemize}
    \item Criterios: AIC, BIC, validación cruzada
    \item Métodos: best subset, forward, backward, stepwise
    \item Siempre combinar con conocimiento del dominio
\end{itemize}

\vspace{0.3cm}
\textbf{Regresión polinomial:}
\begin{itemize}
    \item Para relaciones no lineales
    \item Crear variables $X^2, X^3, \ldots$
    \item Elegir grado con AIC/BIC/CV
    \item Cuidado con sobreajuste (grados altos)
\end{itemize}

\vspace{0.5cm}
\textbf{Próxima clase:} Correlación y ANCOVA
\end{frame}

\end{document}
